\documentclass[letterpaper,11pt]{article}

\usepackage[empty]{fullpage}
\usepackage[hidelinks]{hyperref}
\usepackage{fontawesome5}
\usepackage{enumitem}
\usepackage{parskip}

% Adjust margins
\addtolength{\oddsidemargin}{-0.4in}
\addtolength{\evensidemargin}{-0.4in}
\addtolength{\textwidth}{0.8in}
\addtolength{\topmargin}{-0.5in}
\addtolength{\textheight}{1.0in}

\ifdefined\pdfgentounicode
  \input{glyphtounicode}
  \pdfgentounicode=1
\fi

\begin{document}

\begin{center}
    {\Huge \scshape Juan Becerra} \\ \vspace{4pt}
    \textbf{\large Software Engineer -- Systems, Tools \& Gameplay} \\ \vspace{4pt}
    \small Seattle, WA ~ $|$ ~
    \raisebox{-0.1\height}\faPhone\ (208) 550-9191 ~ $|$ ~
    \href{mailto:becerrajuan007@gmail.com}{\raisebox{-0.2\height}\faEnvelope\ \underline{becerrajuan007@gmail.com}} ~ $|$ ~
    \href{https://hoowan.dev}{\raisebox{-0.2\height}\faGlobe\ \underline{hoowan.dev}} \\ \vspace{2pt}
    \href{https://linkedin.com/in/hoowan/}{\raisebox{-0.2\height}\faLinkedin\ \underline{linkedin.com/in/hoowan}} ~ $|$ ~
    \href{https://github.com/hoowan-dev}{\raisebox{-0.2\height}\faGithub\ \underline{github.com/hoowan-dev}}
\end{center}

\vspace{4pt}
\hrule
\vspace{14pt}

\noindent October 26, 2026

\vspace{8pt}
\noindent \textbf{Yen Feng} \\
Hiring Manager --- CT - Data \& Insights (Enterprise Intelligence Group) \\
Electronic Arts \\
Kirkland, WA

\vspace{8pt}
\noindent \textbf{Re: Software Engineer III - Enterprise Intelligence (Role ID: 215963)}

\vspace{8pt}
\noindent Dear Yen Feng,

I am writing to express my strong interest in the \textbf{Software Engineer III -- Enterprise Intelligence (Role ID: 215963)} position within Electronic Arts' Central Technology Data \& Insights organization in Kirkland. Having spent the past five years at Electronic Arts---first at \textbf{EA Maxis Studios} on \emph{Project Rene} (\emph{The Sims}) and currently at \textbf{EA Lifestyle} here in Kirkland---I have built a strong foundation in high-throughput C++ and Python systems, runtime data pipelines, and cross-studio engineering tools. I am excited by the opportunity to bring my familiarity with EA's live-service ecosystem and my passion for scalable systems engineering to the Enterprise Intelligence Group.

In my current role at \textbf{EA Lifestyle}, I engineered a high-throughput Unreal Engine C++ plugin that extracts and exports runtime character trajectory, pose, and velocity data to populate motion-matching query databases. Alongside data export pipelines, I architected an extensible Python/Maya automation tool adopted across multiple internal EA studios---reducing redundant manual workflows by over 70\%---and integrated PyTorch and ONNX Runtime pipelines for automated asset tagging. Previously at \textbf{EA Maxis Studios}, I architected a multi-threaded, asynchronous path-planning and navigation API in C++ for \emph{Project Rene}, operating within strict frame and latency budgets ($<$1.5ms). Through deep profiling and optimization of compute and memory across concurrent evaluation graphs and query sweeps, I eliminated runtime bottlenecks during dense simulation scenarios while partnering across disciplines to build diagnostic and visual debugging tools.

While my recent roles at EA have focused on runtime gameplay, animation, and cross-studio tooling, the underlying engineering challenges---designing low-latency multi-threaded systems, building reliable data extraction and query pipelines, optimizing compute and memory budgets, and driving developer productivity across studios---align closely with the work of the Enterprise Intelligence team. Combined with my proficiency in C++, Java, Python, and SQL, my familiarity with AWS, Docker, and CI/CD workflows, and my ongoing graduate coursework in Machine Learning \& Deep Learning at the University of Washington, I am eager to apply my systems background to EA's unified Big Data platform.

Thank you for your time and consideration. As an internal EA engineer already based in the Seattle/Kirkland area, I look forward to the opportunity to discuss how my experience can contribute to the Enterprise Intelligence Group.

\vspace{10pt}
\noindent Sincerely,

\vspace{14pt}
\noindent \textbf{Juan Becerra}

\end{document}

